% !TeX root = ../main.tex
\section{Subtyping}

Polymorphism (\S5) let one term have many types. \textbf{Subtyping} asks a different question: when can a term of \emph{one specific} type safely stand in for a different, specific type? The answer becomes a relation $\tau_1 <: \tau_2$ (``$\tau_1$ is a subtype of $\tau_2$''), read informally as ``anywhere a $\tau_2$ is expected, a $\tau_1$ will do''.

\subsection{the subsumption rule}

Subtyping enters the type system through exactly one new rule, \textbf{subsumption}: if you've already shown $e$ has type $\tau_1$, and $\tau_1$ is a subtype of $\tau_2$, you may also conclude $e$ has type $\tau_2$.
\[
\infer[\textsc{T-Sub}]
  {\Gamma \vdash e : \tau_2}
  {\Gamma \vdash e : \tau_1 \quad \tau_1 <: \tau_2}
\]
Everything else about the type system (\textsc{T-Var}, \textsc{T-App}, etc., from \S3) is unchanged; subsumption is a single extra escape hatch you can invoke at any point in a derivation. All the interesting design work in a subtyping system goes into defining $<:$ itself, since that's what determines which programs subsumption actually rescues.

\subsection{record subtyping: width and depth}

The clearest setting to build intuition is records. Suppose we have
\[
  \{\mathit{name}: \mathbf{string},\ \mathit{age}: \tyint,\ \mathit{email}: \mathbf{string}\}
\]
and a function expecting $\{\mathit{name}: \mathbf{string},\ \mathit{age}: \tyint\}$. Intuitively, passing the three-field record should be fine — the function only ever looks at \lstinline|name| and \lstinline|age|, and the extra \lstinline|email| field is simply ignored. This is \textbf{width subtyping}: a record with \emph{more} fields is a subtype of a record with fewer, provided the shared fields match.
\[
  \{\ell_1:\tau_1,\ \dots,\ \ell_n:\tau_n,\ \ell_{n+1}:\tau_{n+1},\ \dots,\ \ell_m:\tau_m\} \;<:\; \{\ell_1:\tau_1,\ \dots,\ \ell_n:\tau_n\}
\]
Separately, \textbf{depth subtyping} says that if the fields' types can \emph{themselves} be replaced by subtypes, the whole record type follows suit:
\[
  \text{if } \tau_i <: \sigma_i \text{ for every } i, \text{ then } \{\ell_1:\tau_1,\dots,\ell_n:\tau_n\} <: \{\ell_1:\sigma_1,\dots,\ell_n:\sigma_n\}
\]
So $\{\mathit{pet}: \{\mathit{name}:\mathbf{string}, \mathit{breed}:\mathbf{string}\}\}$ is a depth-subtype of $\{\mathit{pet}: \{\mathit{name}:\mathbf{string}\}\}$, because the nested \lstinline|pet| field's type is itself narrowed by width subtyping. Width and depth are independent knobs — a full structural record subtyping relation combines both.

\subsection{variance: why functions flip the arrow}

Now the part that trips everyone up at least once: how does subtyping interact with function types? The rule is
\[
\infer[\textsc{S-Arrow}]
  {\tau_1 \to \tau_2 \;<:\; \sigma_1 \to \sigma_2}
  {\sigma_1 <: \tau_1 \quad \tau_2 <: \sigma_2}
\]
Notice the argument types are backwards: to conclude $\tau_1\to\tau_2 <: \sigma_1\to\sigma_2$, we need $\sigma_1 <: \tau_1$ (not $\tau_1<:\sigma_1$). We say function types are \textbf{contravariant} in the argument position and \textbf{covariant} in the return position.

This isn't an arbitrary convention — it's forced by safety. Ask: if I have a function $f : \tau_1\to\tau_2$ and I want to use it \emph{as if} it had type $\sigma_1\to\sigma_2$, what does the caller assume?

\begin{itemize}
  \item The caller, expecting $\sigma_1\to\sigma_2$, is going to call $f$ with an argument of type $\sigma_1$. For that call to be safe, $f$ needs to accept anything of type $\sigma_1$ — i.e., $f$'s actual domain $\tau_1$ must be \emph{at least as general} as $\sigma_1$, so $\sigma_1<:\tau_1$. If $\tau_1$ were narrower than $\sigma_1$ (say $f$ only accepts dogs but the caller might hand it any animal), the call could pass something $f$ doesn't know how to handle. Hence the flip: the wider type is the subtype's domain requirement runs \emph{backwards}.
  \item The caller expects to get back a $\sigma_2$. $f$ actually returns a $\tau_2$. For the caller to be safe treating that result as a $\sigma_2$, we need $\tau_2 <: \sigma_2$ — the actual return value must be usable wherever a $\sigma_2$ was expected. That's the ordinary, forward direction: covariant.
\end{itemize}
In short: it's always safe to accept \emph{more} than you're asked to promise, and to promise \emph{more} than you're asked to accept — never the other way around.

\subsection{Liskov substitution}

This whole idea has a name from object-oriented design: the \textbf{Liskov substitution principle}, which says, informally, that if $S$ is a subtype of $T$, objects of type $T$ in a program should be replaceable with objects of type $S$ \emph{without altering any of the desirable properties of that program} (correctness, not just compiling). That is precisely what subsumption formalizes: $\tau_1<:\tau_2$ is exactly the license to substitute a $\tau_1$-typed thing wherever a $\tau_2$ was expected, and the variance rule above is what keeps that substitution actually sound rather than merely syntactically convenient.

\subsection{nominal vs.\ structural subtyping}

Back in \S1 we flagged nominal vs.\ structural typing as an axis to revisit once we had subtyping in hand — here it is. The question is: what does it take for $\tau_1 <: \tau_2$ (or even $\tau_1 = \tau_2$) to hold — a shared \emph{name} in the program's declarations, or merely a compatible \emph{shape}?

\textbf{Nominal} systems require an explicit declaration of the relationship. Java is the textbook case: a class only satisfies an interface if it says so with \lstinline|implements|.
\begin{lstlisting}[language=Java]
interface Flyer { void fly(); }

class Duck {
    public void fly() { System.out.println("flap flap"); }
}

// error: Duck is not a Flyer -- shape matches, name doesn't
Flyer f = new Duck();
\end{lstlisting}
\lstinline|Duck| has exactly the method \lstinline|Flyer| requires, but the compiler rejects the assignment anyway: without \lstinline|class Duck implements Flyer|, there is no subtyping relationship, full stop, regardless of shape.

\textbf{Structural} systems ask only whether the shape lines up. TypeScript is the clean example:
\begin{lstlisting}
interface Flyer { fly(): void; }

class Duck {
  fly(): void { console.log("flap flap"); }
}

const f: Flyer = new Duck();  // fine: Duck has a fly() method
\end{lstlisting}
\lstinline|Duck| never mentions \lstinline|Flyer| anywhere, and TypeScript accepts it regardless — having the right method with the right signature \emph{is} being a \lstinline|Flyer|. Go's interfaces work the same way (satisfied implicitly, by method set), which is why Go's standard idiom is to define small interfaces and let any type that happens to implement the right methods satisfy them, with zero declared relationship:
\begin{lstlisting}[language=Go]
type Flyer interface {
    Fly()
}

type Duck struct{}
func (d Duck) Fly() { fmt.Println("flap flap") }

var f Flyer = Duck{}  // fine: Duck's method set satisfies Flyer
\end{lstlisting}
Same shape-compatible-but-unrelated scenario as the Java example, opposite outcome — that contrast is the whole nominal/structural distinction in miniature. Nominal typing buys you intentionality (a \lstinline|USD| and an \lstinline|EUR| that both wrap a \lstinline|double| won't accidentally satisfy each other's interfaces); structural typing buys you flexibility (you can satisfy an interface you didn't know existed when you wrote your type, which is exactly what makes duck typing safe with static checking).

\subsection{Exercises}
\begin{itemize}
  \item Using width and depth subtyping, show $\{x:\tyint, y:\tyint, z:\{a:\tyint,b:\tyint\}\} <: \{x:\tyint, z:\{a:\tyint\}\}$, naming which rule (width or depth, and where) licenses each step.
  \item Given $\tau_1\to\tau_2 <: \sigma_1\to\sigma_2$ requires $\sigma_1<:\tau_1$ and $\tau_2<:\sigma_2$, construct a concrete unsafe program that \emph{would} typecheck if the argument-position rule were instead covariant ($\tau_1<:\sigma_1$) rather than contravariant, and explain exactly where it goes wrong at runtime.
  \item Translate the Java \lstinline|Duck|/\lstinline|Flyer| example into Go or TypeScript and confirm it typechecks without any declared relationship. Then translate the reverse: a genuine subtype relationship declared in TypeScript via \lstinline|extends|, and check whether removing the \lstinline|extends| clause (but keeping the same members) still typechecks.
  \item State the Liskov substitution principle in your own words and give an example (real or invented) of a subclass that satisfies a supertype's method signatures but \emph{violates} LSP in spirit (e.g. the classic \texttt{Square extends Rectangle} problem). What goes wrong even though the type checker is happy?
  \item Is subtyping in this section's record calculus reflexive ($\tau<:\tau$) and transitive (if $\tau_1<:\tau_2<:\tau_3$ then $\tau_1<:\tau_3$)? Argue briefly for both properties using the rules given.
\end{itemize}
